\documentclass[conference]{IEEEtran}

\usepackage{graphicx}
\usepackage{amsmath,amssymb}
\usepackage{float}
\usepackage{booktabs}
\usepackage{multirow}

\begin{document}

\title{
    \textbf{Benchmark Report: 520.omnetpp\_r} \\
    \text{Register Spill Analysis on RISC-V via gem5}
}

\author{
    \IEEEauthorblockN{Lütfullah Burak Kaya}
    \IEEEauthorblockA{
        Ozyegin University \\
        burak.kaya.50711@ozu.edu.tr
    }
    \and
    \IEEEauthorblockN{Ismail Akturk}
    \IEEEauthorblockA{
        Ozyegin University \\
        ismail.akturk@ozyegin.edu.tr
    }
}

\newcommand{\LongDate}{%
    \number\day\space
    \ifcase\month
    \or January\or February\or March\or April\or May\or June%
    \or July\or August\or September\or October\or November\or December%
    \fi
    \space \number\year
}

\twocolumn[
\begin{@twocolumnfalse}
\maketitle
\begin{center}{\small \LongDate}\end{center}
\vspace{1em}
\end{@twocolumnfalse}
]

% ================================================

\begin{abstract}
This report presents the register spill characterization of the
\texttt{520.omnetpp\_r} benchmark from SPEC CPU2017, executed on a
RISC-V (RV64GC) target using the gem5 TimingSimpleCPU simulator.
A custom SpillDetector was integrated into gem5 to count store-then-load
patterns on the stack within a user-defined Region of Interest (ROI).
The test dataset (\texttt{General} configuration, run~0,
\texttt{sim-time-limit=0.003s}) was used as the input. Results show
a spill rate of \textbf{7.20\%} within the ROI --- fourth highest
across all SPEC CPU2017 benchmarks evaluated in this study ---
corresponding to 895 million detected spill events across 12.4 billion
ROI instructions. The heavily object-oriented OMNeT++ event-dispatch
loop maintains a large set of simultaneously live module pointers,
message objects, and scheduling state, creating persistent register
pressure throughout the network simulation.
\end{abstract}

% =========================================================
\section{Benchmark Overview}

\texttt{520.omnetpp\_r} is the SPEC CPU2017 rate variant of
OMNeT++~4.0, a discrete-event network simulator. The benchmark
simulates a large Ethernet LAN (\texttt{LargeNet}) using the INET
framework, exercising event scheduling, message passing between
network modules, MAC layer processing, and packet forwarding.
OMNeT++ is one of the most object-oriented workloads in the SPEC
suite: the event loop dispatches messages through a deep C++ class
hierarchy of simulated network components via virtual function calls,
maintaining complex module state and scheduling bookkeeping.

The RISC-V binary (\texttt{omnetpp\_r.riscv}) was compiled with gem5
ROI markers inserted around the \texttt{setupUserInterface()} call
in \texttt{src/simulator/main.cc}:

\begin{itemize}
    \item \texttt{m5\_work\_begin(0,0)} --- before \texttt{setupUserInterface()}
    \item \texttt{m5\_work\_end(0,0)}   --- after  \texttt{setupUserInterface()}
\end{itemize}

\texttt{setupUserInterface()} instantiates the Cmdenv command-line
environment, which builds the simulated network, initialises all
modules, runs the event loop for the configured simulation time,
and calls \texttt{finish()} on each module. This call encompasses
the entire simulated network execution.

\subsection{Input Datasets}

SPEC CPU2017 provides three dataset sizes for \texttt{omnetpp\_r},
differing in simulated network time and topology scale
(Table~\ref{tab:inputs}).

\begin{table}[H]
\centering
\caption{omnetpp\_r Input Dataset Sizes}
\label{tab:inputs}
\begin{tabular}{lll}
\toprule
\textbf{Dataset} & \textbf{Sim time limit} & \textbf{Description} \\
\midrule
test    & 0.003\,s & LargeNet Ethernet, short burst \\
train   & 0.1\,s   & LargeNet Ethernet, medium run \\
refrate & 2.25\,s  & LargeNet Ethernet, full run \\
\bottomrule
\end{tabular}
\end{table}

The \textbf{test} dataset was selected for this study. Despite the
short simulated time of 3\,ms, the OMNeT++ event loop generates
12.4 billion RISC-V instructions because each discrete network event
requires traversing the full C++ object hierarchy: scheduler lookup,
message delivery through module virtual functions, MAC layer state
machine evaluation, and statistics recording. The network topology
(\texttt{LargeNet}) contains numerous Ethernet switches, hosts, and
buses, each modelled as a separate C++ object.

% =========================================================
\section{Simulation Setup}

\subsection{gem5 Configuration}

\begin{itemize}
    \item \textbf{CPU model:}   TimingSimpleCPU
    \item \textbf{ISA:}         RISC-V RV64GC
    \item \textbf{Caches:}      enabled (default L1 I/D + L2)
    \item \textbf{Memory:}      4\,GB simulated physical RAM
    \item \textbf{Spill mode:}  summary only (\texttt{spill\_verbose=False})
\end{itemize}

The full gem5 invocation was:

\begin{verbatim}
./build/RISCV/gem5.opt
  --outdir=benchmarks/omnetpp/m5out_test
  configs/deprecated/example/se.py
  --cmd=benchmarks/omnetpp/omnetpp_r.riscv
  --options="-c General -r 0"
  --cpu-type=TimingSimpleCPU
  --caches
  --mem-size=4GB
\end{verbatim}

The binary locates \texttt{omnetpp.ini} in the gem5 working directory
via a symbolic link, with the \texttt{ned-path} pointing to the
benchmark's \texttt{ned/} directory containing the LargeNet NED
module definitions.

% =========================================================
\section{Simulation Results}

The simulation completed successfully. OMNeT++ ran the
\texttt{General} configuration for run~0, processed all network
events up to the 3\,ms simulation time limit, called
\texttt{finish()} on all modules, and exited cleanly.

\subsection{Pre-ROI Context}

Before \texttt{setupUserInterface()} was entered, the binary
executed the OMNeT++ runtime initialisation: dynamic library
resolution, global object construction, NED file loading (13 files),
and configuration parsing. Table~\ref{tab:preroi} shows the global
counters at ROI entry.

\begin{table}[H]
\centering
\caption{Global Counters at ROI Entry}
\label{tab:preroi}
\begin{tabular}{lr}
\toprule
\textbf{Metric} & \textbf{Value} \\
\midrule
Total instructions (pre-ROI) & 1,623,908 \\
Total spills detected (pre-ROI) & 56,740 \\
Pre-ROI spill rate & 3.49\% \\
\bottomrule
\end{tabular}
\end{table}

\subsection{ROI Spill Statistics}

Table~\ref{tab:roi} presents the spill statistics collected within
the ROI.

\begin{table}[H]
\centering
\caption{ROI Spill Statistics --- 520.omnetpp\_r (test, 0.003\,s)}
\label{tab:roi}
\begin{tabular}{lr}
\toprule
\textbf{Metric} & \textbf{Value} \\
\midrule
ROI instructions         & 12,426,280,892 \\
ROI stores               &  1,491,466,063 \\
ROI loads                &  2,793,818,477 \\
\midrule
\textbf{ROI spills}      & \textbf{894,619,681} \\
\textbf{Spill rate}      & \textbf{7.1994\%} \\
\bottomrule
\end{tabular}
\end{table}

\subsection{Timing}

\begin{table}[H]
\centering
\caption{Simulation Timing}
\label{tab:timing}
\begin{tabular}{lr}
\toprule
\textbf{Metric} & \textbf{Value} \\
\midrule
Simulated time           & 30.17\,s \\
Host wall time           & 7{,}777.80\,s ($\approx$2.16\,hours) \\
Total simulated instructions & 12,429,400,619 \\
\bottomrule
\end{tabular}
\end{table}

% =========================================================
\section{Analysis}

\subsection{Spill Rate Interpretation}

The measured spill rate of \textbf{7.20\%} places
\texttt{omnetpp\_r} fourth among all evaluated SPEC CPU2017
benchmarks (Table~\ref{tab:compare}), between
\texttt{500.perlbench\_r} (7.28\%) and \texttt{502.gcc\_r} (6.92\%).

\begin{table}[H]
\centering
\caption{Cross-Benchmark Spill Rate Comparison}
\label{tab:compare}
\begin{tabular}{lrr}
\toprule
\textbf{Benchmark} & \textbf{ROI Spills} & \textbf{Spill Rate} \\
\midrule
507.cactuBSSN\_r & 3,782,099,438 & 7.8756\% \\
511.povray\_r    &   174,359,922 & 7.7809\% \\
500.perlbench\_r &     1,070,055 & 7.2838\% \\
520.omnetpp\_r   &   894,619,681 & \textbf{7.1994\%} \\
502.gcc\_r       &     1,108,838 & 6.9185\% \\
526.blender\_r   &    61,845,869 & 6.2193\% \\
523.xalancbmk\_r &    18,381,250 & 4.6427\% \\
541.leela\_r     & 1,158,236,271 & 4.4438\% \\
519.lbm\_r       &   274,130,440 & 4.2130\% \\
531.deepsjeng\_r & 1,928,903,656 & 4.1749\% \\
525.x264\_r      &   767,981,981 & 3.2812\% \\
544.nab\_r       &   151,310,616 & 2.2612\% \\
538.imagick\_r   &     2,424,899 & 2.0601\% \\
557.xz\_r        &    38,411,428 & 1.9069\% \\
505.mcf\_r       &   446,113,136 & 1.8263\% \\
508.namd\_r      &   327,718,955 & 1.0269\% \\
\bottomrule
\end{tabular}
\end{table}

\subsection{Store vs.\ Load Balance}

The load count (2.79B) is $1.87\times$ higher than the store count
(1.49B). This load-dominated ratio characterises OMNeT++'s event
dispatch pattern: each event requires reading the module's state
(current queue, pending timers, MAC address tables) and reading
message header fields to determine how to process and forward the
packet. Writes occur when updating module state, recording statistics,
and scheduling new events. The event loop reads far more than it writes
because evaluating a network event involves traversing multiple
linked data structures before producing a single scheduling write.

\subsection{Spill Fraction of Loads}

Within the ROI, the ratio of spills to total loads is:
\[
\frac{894{,}619{,}681}{2{,}793{,}818{,}477} \approx 32.0\%
\]
Approximately 1 in 3 load instructions is a spill reload. This is
the second highest spill-to-load fraction observed in the study,
behind only \texttt{502.gcc\_r} (37.5\%). The OMNeT++ event loop
must simultaneously maintain: a pointer to the current event, the
receiving module pointer, the message object pointer, the simulation
clock, the future event set iterator, and the per-module state
structure. The deep C++ virtual dispatch chain --- from the scheduler
to the module's \texttt{handleMessage()} to the specific protocol
handler --- spills and reloads these pointers at every level,
as RISC-V's 32 registers cannot hold all of them live across
nested virtual calls.

\subsection{Pre-ROI Spill Density}

The pre-ROI spill rate of 3.49\% ($56{,}740$ spills over
$1{,}623{,}908$ instructions) reflects the OMNeT++ runtime
initialisation: loading NED module definitions, parsing the ini
file, and registering global C++ objects. This moderate pre-ROI
rate indicates that even the OMNeT++ startup phase involves
non-trivial register pressure from its object registration framework.

\subsection{ROI Coverage}

The ROI accounts for $12{,}426{,}280{,}892$ out of $12{,}429{,}400{,}619$
total simulated instructions:
\[
\frac{12{,}426{,}280{,}892}{12{,}429{,}400{,}619} \approx 99.97\%
\]
The near-perfect ROI coverage confirms that
\texttt{setupUserInterface()} encompasses virtually all benchmark
computation. The 0.03\% outside the ROI corresponds to
pre-ROI startup and post-simulation shutdown.

\subsection{Discrete-Event Simulation Register Pressure}

The 7.20\% spill rate is structurally driven by OMNeT++'s
object-oriented event dispatch architecture. Each call to
\texttt{cSimulation::doOneEvent()} invokes a chain of virtual
functions that must propagate the current message, the receiving
module, and the simulation context through multiple layers of
abstraction:
\begin{itemize}
    \item \texttt{cScheduler} $\rightarrow$ \texttt{cSimpleModule::handleMessage()}
    \item MAC layer dispatch (EtherMAC, EtherHub, EtherSwitch)
    \item Statistics recording (\texttt{cOutVector}, \texttt{cStdDev})
\end{itemize}
At each level, the caller's pointers must be preserved across the
virtual call, and RISC-V's 32-register file cannot accommodate all
live pointers simultaneously. The result is that every network event
produces several spill-reload cycles in the dispatch chain, yielding
the observed 32.0\% spill-to-load fraction.

% =========================================================
\section{Conclusion}

The \texttt{520.omnetpp\_r} benchmark exhibits a register spill rate
of 7.20\% within the \texttt{setupUserInterface()} ROI on RISC-V
RV64GC, the fourth highest across all 16 evaluated SPEC CPU2017
benchmarks. With 895 million spill events across 12.4 billion ROI
instructions and a 99.97\% ROI coverage, the benchmark is
comprehensively characterised by this run. The 1.87:1 load-to-store
ratio and 32.0\% spill-to-load fraction reflect the read-dominated
event dispatch pattern of the OMNeT++ simulator, where each network
event triggers a deep virtual call chain that continuously exhausts
RISC-V's register file. These results confirm that object-oriented
discrete-event simulation is among the most register-pressure
intensive workload classes on RISC-V RV64GC.

\end{document}
